\chapter{State of the art}
\label{ch:sota}

This thesis sits at the intersection of four literatures rarely brought together in one system: teleoperation of remote manipulators, where an operator commands a robot arm through a haptic or kinematic interface; shared control and shared autonomy, where an inferred estimate of operator intent contributes to the command actually executed; safety-critical control, where control barrier functions (CBFs) realized inside a quadratic program (QP) keep the commanded motion inside a collision-free set by formal guarantee rather than heuristic penalty; and perception for safety-critical control, where the obstacle geometry that safety layer relies on is measured by the robot's own camera rather than given to it beforehand. Each field is mature on its own. What this chapter argues, and what Section~\ref{sec:sota-positioning} makes explicit against the closest prior system, is that their intersection --- safety enforced strictly after arbitration, on real hardware, with the safety filter's own dual variables closing the loop back to both controller and operator --- remains sparsely populated. Section~\ref{sec:sota-teleoperation} reviews bilateral teleoperation and the motion mappings this thesis treats as a factor of its user study. Section~\ref{sec:sota-shared-control} surveys shared autonomy from intent inference to arbitration and isolates the failure mode of linear blending that motivates the rest of the chapter. Section~\ref{sec:sota-haptic-fixtures} covers haptic shared control as a mechanism for making assistance felt rather than only computed. Section~\ref{sec:sota-safety-critical} is the technical core, tracing control barrier functions from potential fields to unified control-Lyapunov-function--control-barrier-function quadratic programs and their use inside the teleoperation loop. Section~\ref{sec:sota-perception} reviews the perception pipelines that build the obstacle model such a controller consumes. Section~\ref{sec:sota-positioning} closes the chapter by positioning this thesis's contribution against the reviewed work.

\section{Teleoperation of remote manipulators}
\label{sec:sota-teleoperation}

Bilateral teleoperation closes the loop between operator and remote manipulator in both directions: the operator's motion drives the robot, and a force computed from the robot's interaction with its environment is rendered back to the operator's hand. Hokayem and Spong survey a century's worth of architectures for this loop, from simple position--position and force--position schemes to four-channel architectures that pass both position and force information in each direction~\cite{hokayem2006bilateral}. The central difficulty, established in Lawrence's two-port formulation, is that transparency (the rendered impedance matching the true remote impedance) and stability trade against each other whenever the channel introduces delay or the haptic display has finite bandwidth~\cite{lawrence1993stability}.

\label{sec:sota-mappings}
\begin{sloppypar}
A teleoperation interface maps the operator's hand motion, bounded by the haptic device's own workspace, onto the larger and differently shaped workspace of the remote manipulator. Two mappings dominate the literature, and Kim, Tendick, Ellis, and Stark's comparison remains the reference study establishing that they are not interchangeable~\cite{kim1987comparison}. In position control, the handle pose maps directly to a desired remote pose; because the device workspace is smaller than the manipulator's reach, position control requires indexing, also called clutching --- the operator periodically disengages the mapping, repositions the handle, and re-engages --- trading continuous correspondence for an added motor step at every indexing event. In rate control, handle displacement from a centred rest pose maps to a commanded remote velocity, removing the workspace-mismatch problem at the cost of an integrating hand-to-end-effector relationship generally found less intuitive for fine positioning. Kim et al. found rate control to dominate once the manipulator's dynamics differ substantially from the hand's own, and position control to dominate when the two are kinematically similar and indexing cost is low~\cite{kim1987comparison}. The distinction reaches deeper than ergonomics: zero handle displacement means something different in each mapping. In rate control it commands zero remote velocity, a natural rest state; in position control with indexing it can mean the handle is disengaged and moving freely while the remote arm holds still, so the same physical handle state carries opposite consequences for what the robot is doing. This asymmetry is why this thesis treats the choice of motion mapping as a factor of the user study rather than a fixed implementation decision.
\end{sloppypar}

Both mappings face a further problem specific to the wrist: a device's rotational workspace is typically smaller and differently shaped than the manipulator's, so the operator's and the end-effector's orientations drift apart unless something keeps them linked. Allenspach, Lawrance, Tognon, and Siegwart's rate-control recentering scheme for aerial teleoperation returns the master toward a fixed rest pose whenever the operator releases the handle, so the vehicle settles at whatever pose it held at release~\cite{allenspach2022bilateral}; the rest pose itself is static, not recomputed from the vehicle's own orientation, so recentering aligns the two only passively, after the fact of release, rather than continuously. W\"ohle and Gebhard's head- and eye-gaze interface for Cartesian robot control samples the current end-effector orientation as a zero reference at the moment control switches into an orientation mode, then maps subsequent input as a difference from that one sample~\cite{wohle2021gaze}; the reference is refreshed only at that single transition, and no restoring force is rendered to pull the input device back toward it. Neither continuously re-centers a rate-control rest pose on the remote's live orientation, nor renders an active torque that pulls the input device's orientation back into alignment during indexing --- the two synchronisation mechanisms this thesis pairs with its own rate- and position-control mappings.

\section{Shared control and shared autonomy}
\label{sec:sota-shared-control}

\subsection{Taxonomies and levels of autonomy}
\label{sec:sota-taxonomies}

Physical human--robot interaction spans a spectrum from pure teleoperation, in which every degree of freedom is commanded by the operator, to full autonomy, in which the robot acts without operator input, and the literature draws its terminology from where a system sits on that spectrum. Selvaggio, Cognetti, Nikolaidis, Ivaldi, and Siciliano's survey organizes physical human--robot interaction along who commands, who senses, and how authority is negotiated, separating \emph{shared control} --- continuous, low-level co-regulation of a single command signal --- from \emph{shared autonomy}, in which a discrete estimate of the operator's task-level goal informs an autonomous policy combined with the operator's own command~\cite{selvaggio2021autonomy}. Losey, McDonald, Battaglia, and O'Malley's review decomposes any such system into three stages that recur across this chapter: intent detection (Section~\ref{sec:sota-intent}), arbitration of authority (Section~\ref{sec:sota-arbitration}), and communication of the system's state back to the operator (Section~\ref{sec:sota-haptic-fixtures})~\cite{losey2018review}. Both surveys agree that arbitration is where the design space is largest and least standardized, borne out by the diversity of signals reviewed below, from goal-probability confidence to distance, entropy, and directional agreement.

\subsection{Inferring operator intent}
\label{sec:sota-intent}

Formal intent inference in robotics traces back to inverse decision-theoretic models developed outside robotics: Baker, Saxe, and Tenenbaum model an observed agent's actions as approximately optimal with respect to an unknown goal and invert a Bayesian model of planning to recover a posterior over goals from partial trajectories~\cite{baker2009action}, and Ziebart, Maas, Bagnell, and Dey's maximum-entropy formulation of inverse reinforcement learning gives the same noisily rational assumption a principled probabilistic form, modelling an observed action's likelihood as exponential in its cost under each candidate goal~\cite{ziebart2008maxent}. Javdani, Admoni, Pellegrinelli, Srinivasa, and Bagnell cast shared autonomy as a partially observable Markov decision process (POMDP) in which the unobserved goal is a hidden state, approximating the hindsight-optimization policy in real time by solving fully observed problems conditioned on each candidate goal~\cite{javdani2018hindsight}. Jain and Argall instantiate the same observation model inside a recursive Bayesian filter that updates a discrete belief over goals online from the operator's own commanded action, a computationally light alternative well suited to a control loop running at haptic rates~\cite{jain2020probabilistic}. Muelling et al. extend the same infrastructure to a brain--computer interface, where the observed action is itself extremely noisy~\cite{muelling2017autonomyinfused}, and Song, Li, Aertbelien, and Detry replace the discrete-goal belief with a learned trajectory predictor, arbitrating from a continuous forecast of where the operator's motion is heading~\cite{song2024trajectron}. A structural hazard runs through the recursive-filter family regardless of observation model: if the observation conditioned on is derived from the robot's own executed motion rather than the operator's raw command, the estimator partially observes its own output and belief can become self-reinforcing --- a goal that momentarily attracts more authority pulls the end-effector toward it, read back as further evidence for that goal, biasing the filter toward stickiness regardless of a genuine change of mind. Conditioning strictly on the operator's own commanded action, before any autonomous contribution is added, avoids closing this loop.

\subsection{Arbitrating authority}
\label{sec:sota-arbitration}

Once a belief over goals exists, a second decision governs how much authority the autonomous policy receives at any instant. Dragan and Srinivasa formalize assistive teleoperation as a convex combination of the user's raw input and an autonomous policy, arbitrated by a weight that is a function of belief confidence~\cite{dragan2012formalizing}; their policy-blending formalism generalizes this so the weight can also depend on distance to the goal or on belief entropy~\cite{dragan2013policy}. Nikolaidis, Hsu, and Srinivasa extend arbitration into a mutual-adaptation setting, modelling the human as also adapting to the robot's revealed policy, so authority allocation becomes a coupled estimation problem rather than a one-shot computation~\cite{nikolaidis2017mutualadaptation}. All of these designs key the arbitration weight to a scalar summary of confidence. Oh, Toussaint, and Mainprice depart from this family by arbitrating instead on the \emph{disagreement} between candidate sub-policies --- an operator command is trusted more when independently reasonable strategies converge on similar actions and less when they diverge, regardless of how confident any single strategy is in isolation~\cite{oh2021disagreement}. A confidence-based weight can be high while the confident policy points where the operator is actively resisting, whereas a disagreement-based weight is low exactly when human and autonomous intent conflict. This directional framing is the closest prior formulation to arbitrating on the alignment between the operator's commanded twist and the policy's twist, taken up again in Section~\ref{sec:sota-positioning}. A further concern is independent of how the weight is computed and applies to the motion the blend produces: Dragan, Lee, and Srinivasa formalise \emph{legible} motion --- shaped so that an observer infers its goal quickly and confidently from a partial trajectory --- as distinct from merely \emph{predictable} motion, and show in a user study that the two objectives can be in direct conflict~\cite{dragan2013legibility}.

\label{sec:sota-linear-blending-limits}
However the arbitration weight is computed, the blended command it produces is not itself guaranteed to respect every safety constraint. Chaubey, Verdoja, Deka, and Kyrki show that once the admissible set is nonconvex, as it generically is around real obstacle geometry, a human command and a robot command can each individually satisfy every constraint while their blend at some intermediate weight violates one, and propose a shared-control law that preserves safety under exactly this nonconvexity~\cite{chaubey2026minimalintervention}. Guler, Pompetzki, Sun, Manschitz, and Peters supply the complementary empirical demonstration inside a full shared-autonomy pipeline: comparing an unconstrained blend against the same blend followed by a control-barrier-function safety filter applied at the inverse kinematics (IK) layer, they show the filtered variant reduces constraint-violation time by more than an order of magnitude and yields a substantially less negative minimum clearance in cluttered scenes, while a soft, cost-based penalty on the same blended command closes only part of that gap~\cite{guler2026barrierik}. Both results converge on the same conclusion: safety cannot be delegated to the arbitration law, nor adequately approximated by a soft penalty on the blend objective --- it requires a hard constraint applied \emph{after} arbitration, on the already-blended command, so the command reaching the robot is checked once, at the last stage, against the true admissible set. This is the hinge on which the rest of this chapter turns: Section~\ref{sec:sota-safety-critical} reviews the machinery providing that hard, post-blend guarantee, and Section~\ref{sec:sota-positioning} returns to this pair of results to situate how this thesis realizes the same principle, with the safety filter validated on physical hardware and the resulting arbitration evaluated systematically in a matched simulation environment.

\subsection{Learned and end-to-end alternatives}
\label{sec:sota-learned}

An alternative to the belief-plus-arbitration pipeline replaces explicit goal inference with a policy learned end to end. Reddy, Dragan, and Levine train a deep reinforcement-learning policy that consumes the raw operator input directly, without forming an explicit belief over a discrete goal set, and show it can outperform a hand-designed confidence-arbitrated blend on tasks where the goal set is difficult to enumerate in advance~\cite{reddy2018sharedrl}. Jeon, Losey, and Sadigh instead learn a low-dimensional latent action space from demonstrations, with the decoder trained so a small, easy-to-produce human input expands into a full, task-appropriate high-dimensional command~\cite{jeon2020latentactions}. Both approaches trade the interpretability of an explicit goal set and arbitration law --- the two design axes this thesis manipulates directly --- for a policy whose reasoning is opaque and whose safety properties are not analysed in closed form, precisely the gap a downstream hard safety constraint is positioned to close.

\subsection{Bimanual and assistive manipulation}
\label{sec:sota-bimanual}

Extending shared control to a bimanual manipulator adds a coordination problem on top of arbitration: two arms may need independent authority allocation while remaining kinematically or task-coupled. Rakita, Mutlu, Gleicher, and Hiatt's shared-control framework lets an operator drive one arm directly while an autonomous controller manages the coordinated motion of the other, validated on physical dual-arm hardware performing handover and joint manipulation~\cite{rakita2019bimanual}. Laghi et al. pursue the same coordination problem from the haptic side, mapping a single-handed operator command to a two-armed humanoid through an interface that preserves the ability to intervene on either arm independently~\cite{laghi2018bimanual}. Manschitz, Guler, Ma, and Ruiken address a complementary problem specific to assisted grasping: sampling candidate grasp poses and checking them against a collision model in real time, so the autonomous side of a shared-autonomy pipeline always proposes a reachable, collision-aware grasp~\cite{manschitz2025samplinggrasp}. None of the three treats the two arms as sharing a single global safety constraint set, leaving open how a bimanual safety filter should behave when one arm's avoidance motion has consequences for the other. Neither Manschitz et al. nor Guler et al. addresses what becomes of the barrier once a grasp actually succeeds: Manschitz et al. withhold their own collision-aware assistance for the entire held-object phase, naming it explicitly as future work, and Guler et al.'s safety formulation carries no notion of a graspable object at all, so nothing in their pick-and-place evaluation of BarrierIK distinguishes an obstacle to avoid from something the operator has just picked up~\cite{guler2026barrierik}.

\section{Haptic assistance in shared control}
\label{sec:sota-haptic-fixtures}

Haptic feedback can render an autonomy's contribution as a force superposed on the operator's own control input, rather than only as a change to the command computed downstream. A recurring principle, articulated by Abbink, Mulder, and Boer, is that this kind of assistance should \emph{bias}, never replace, the operator's own control: the assistive force shifts the effort required to deviate from the suggested behaviour without removing the operator's authority, preserving a continuously variable handoff of authority rather than a discrete mode switch~\cite{abbink2012hapticshared}. Griffiths and Gillespie provide early empirical support in a shared steering task, showing a haptic guidance force improves primary-task tracking without degrading a concurrent secondary task, evidence the assistance is additive rather than substitutive~\cite{griffiths2005sharing}; Mulder, Abbink, and Boer find the same continuously blended authority transition preferred over an abrupt one in a driving handover~\cite{mulder2012sharingcontrol}. Closer to Section~\ref{sec:sota-shared-control}, Zhang, Tron, and Khurshid show that haptic feedback of the autonomy's intervention measurably improves both human--robot agreement and operator satisfaction relative to the same policy rendered silently~\cite{zhang2021hapticagreement} --- evidence that Losey et al.'s communication stage changes how arbitration is experienced and accepted. This work establishes felt assistance as a distinct design axis, one this thesis extends by rendering a safety constraint, rather than a guidance cue, as a haptic force.

\section{Safety-critical control}
\label{sec:sota-safety-critical}

\subsection{From potential fields to formal invariance}
\label{sec:sota-potential-fields}

The earliest reactive approach to manipulator obstacle avoidance treats the environment as a source of artificial forces: Khatib's real-time obstacle-avoidance scheme computes a repulsive potential from each nearby obstacle and adds its negative gradient to the manipulator's task-space control law, producing collision-avoiding motion within a control loop fast enough for online use~\cite{khatib1986}. Maciejewski and Klein extend the idea to kinematically redundant manipulators, projecting the avoidance gradient into the null space of the primary task Jacobian so avoidance and tracking can be pursued simultaneously whenever spare degrees of freedom exist~\cite{maciejewski1985obstacle}. Both methods share a structural weakness: the resulting motion is only as safe as the potential's gradient happens to be at each instant, with no guarantee that the safe region is forward-invariant, that is, that a trajectory starting inside it never leaves. Singletary et al. give this weakness a precise characterization, proving that an artificial potential field is recoverable as a special case of a control barrier function under a restrictive choice of class-$\mathcal{K}$ function, and that outside that special case a potential field generally fails to certify forward invariance where a control barrier function still can~\cite{singletary2021comparative}. Their comparison motivates the shift, taken up next, to control barrier functions as a formal safety certificate.

\subsection{Control barrier functions}
\label{sec:sota-cbf}

A control barrier function formalizes the forward-invariance property a potential field can only approximate. The safe region is described as the set on which a continuously differentiable scalar function of the state stays non-negative, and the barrier condition constrains how fast that function may decrease: its admissible rate of decrease is bounded by a class-$\mathcal{K}$ envelope of its own current value, so a state deep inside the safe set may approach the boundary quickly, while the admissible approach rate is driven to zero as the boundary is reached. Ames, Xu, Grizzle, and Tabuada establish that any Lipschitz-continuous controller meeting this condition pointwise renders the safe set forward-invariant, and that for a control-affine system the inputs meeting it at a given state form an affine half-space --- which is what turns safety into a single linear constraint inside a real-time optimization, rather than a property to be verified offline~\cite{ames2017cbf}. Ames, Coogan, Egerstedt, Notomista, Sreenath, and Tabuada's tutorial survey consolidates this theory and its applications, from legged locomotion to multi-robot coordination~\cite{ames2019ecc}. The condition in this first-order form presupposes that the barrier function responds to the control input at its first derivative; most physically meaningful safety functions, such as a distance-based barrier commanded through acceleration or torque, do not, and Nguyen and Sreenath's exponential control barrier functions extend the condition to that case through a chain of derivative-based conditions parametrized by pole placement~\cite{nguyen2016exponential}, later generalized to arbitrary relative degree by Xiao and Belta~\cite{xiao2022highorder}.

\label{sec:sota-teleop-barriers}
\begin{sloppypar}
Beyond these general guarantees, applying a control-barrier-function safety filter specifically to a teleoperated system raises a question the discussion above does not address: what happens to the safety margin as felt information, not only as a silently enforced constraint. Xu and Sreenath apply a CBF-QP to teleoperated aerial vehicles, filtering the operator's raw velocity command before it reaches the vehicle, with the filter acting purely as a silent safeguard~\cite{xu2018safeteleoperation}. Zhang, Yang, and Khurshid close that gap for the same class of vehicle: their scheme renders the deviation between the commanded and the CBF-filtered velocity as a force on the operator's hand, so the pilot feels the filter intervening rather than only observing its consequence in the vehicle's motion~\cite{zhang2020hapticuav}. Qin, Yi, Fan, and Zhao build a comparable haptic shared-control framework with an interaction-force constraint expressed as a CBF, extending the same principle to contact-rich teleoperation~\cite{qin2025hapticcbf}. Raphalen, Ferro, Amouroux, Robuffo Giordano, and Pacchierotti give this line of work its most direct precedent to the present thesis: rather than deriving a haptic force from the raw command deviation, they render the CBF-QP's own Lagrange multipliers, the same shadow-price dual variables of Section~\ref{sec:sota-clf-cbf-qp}, directly as multi-sensory feedback, reasoning that the multiplier already measures how hard a constraint is fighting the task~\cite{raphalen2026feeling}. This is, on the evidence reviewed here, the closest published instance of using a QP-based safety filter's dual variables as the operator feedback signal; Section~\ref{sec:sota-positioning} returns to it when situating this thesis's own use of the same shadow price for an adaptive constraint-scheduling role inside the controller itself, rather than as external feedback.
\end{sloppypar}

\subsection{Unified CLF--CBF QPs and spurious equilibria}
\label{sec:sota-clf-cbf-qp}

Task tracking and safety pull a real-time controller in opposite directions whenever they conflict, and the standard reconciliation encodes tracking as a soft, relaxable objective and safety as a hard constraint inside one QP solved every tick. A control Lyapunov function certifies convergence to a task goal through a condition of the same algebraic form as the barrier condition but reversed in sign, demanding that a positive-definite measure of task error decay at least at a prescribed rate; admitting a non-negative relaxation variable on that demand, and minimizing it alongside control effort, lets the controller trade convergence rate for safety exactly when the two are in tension, rather than pre-committing to a fixed weighting~\cite{ames2017cbf}. The choice of Lyapunov candidate is not restricted to the quadratic forms used almost universally in this literature: taking the error norm itself is admissible under the nonsmooth control-Lyapunov theory of Clarke, Ledyaev, Sontag, and Subbotin, who prove that asymptotic controllability is equivalent to the existence of a control Lyapunov function that is in general merely continuous, and not differentiable at the origin~\cite{clarke1997asymptotic}. The substitution changes the geometry of the resulting constraint row and not only its algebra: a quadratic candidate yields a gradient that scales with the error and a convergence demand quadratic in it, whereas the norm candidate is positively homogeneous of degree one, so its gradient is the unit error direction and the demanded closing speed grows only linearly --- leaving the row equally well conditioned at metre-scale and at millimetric error. Bhat and Bernstein situate the same candidate in a wider family, showing that a decay exponent below one on such a homogeneous candidate converges in finite rather than merely asymptotic time~\cite{bhat2000finitetime}. The resulting unified CLF--CBF QP is a strictly convex quadratic program, solvable in real time with dual active-set methods such as Goldfarb and Idnani's~\cite{goldfarb1983}, and its Karush--Kuhn--Tucker (KKT) dual variables have the standard convex-optimization interpretation as shadow prices, quantifying the marginal cost of tightening each constraint by one unit~\cite{boyd2004}. The formulation is not free of failure modes: Reis, Aguiar, and Tabuada show that the unified QP can introduce spurious, asymptotically stable equilibria in the closed loop that exist in neither the CLF-only nor the CBF-only vector field alone, arising from the interaction between the two constraints, and that these can trap the system away from its goal even when the goal remains reachable~\cite{reis2021undesirable}. Their result cautions against treating the CLF--CBF QP as solved: real-time feasibility and a formal safety guarantee do not rule out degraded closed-loop behaviour that only a case-by-case analysis can catch.

\subsection{Aggregating many constraints}
\label{sec:sota-aggregation}

A manipulator moving near cluttered geometry is rarely subject to a single barrier constraint: at any instant, many link--obstacle pairs may be close enough to matter, and the active pair can switch discontinuously as the manipulator moves. Glotfelter, Cortes, and Egerstedt address this by composing individual barrier functions with Boolean logic realized through the hard minimum and maximum operators, so a conjunction of constraints becomes a single nonsmooth barrier taking the least of them; they extend CBF forward-invariance theory to this nonsmooth setting, at the cost of a barrier gradient that is undefined, and a resulting QP row that jumps, exactly when the active pair switches~\cite{glotfelter2017}. Rabiee and Hoagg replace the hard minimum with a soft, differentiable surrogate built from the log-sum-exp function, governed by a temperature that controls how closely the surrogate approaches the true minimum, and derive the corresponding smooth barrier condition together with an explicit bound relating the approximation's conservativeness to that temperature, so the aggregate always underestimates the true minimum clearance by a quantifiable margin~\cite{rabiee2025softmin}. Because the soft aggregate is differentiable everywhere, the resulting QP row varies continuously as the closest pair changes, trading bounded conservativeness for the elimination of Glotfelter et al.'s discontinuity. Morton and Pavone show that an aggregated operational-space CBF-QP of this kind scales to hundreds of simultaneous constraints on a full-body system without becoming the computational bottleneck, provided the gradients are assembled efficiently~\cite{morton2025oscbf}, evidence that soft aggregation is a prerequisite for real-time operation once the constraint count grows past a handful of pairs. The aggregate's smoothness is easily mistaken for linearity: even when every individual barrier is a simple, near-affine distance function, the log-sum-exp surrogate is a nonlinear, if convex, function of the state by construction. The resulting per-tick QP row remains a single row rather than one per pair, because the aggregate's gradient is a softmax-weighted blend of the individual gradients, not their conjunction: aggregation trades many simultaneous half-space constraints for one direction that best represents all of them at once. When two obstacles are comparably close and pull in genuinely different directions, no choice of temperature recovers what keeping every pair as its own row would give; only separate rows can enforce independent half-space constraints within the same solve.

\subsection{Reactive optimization for redundant manipulators}
\label{sec:sota-optimization-ik}

Redundant manipulators, with more degrees of freedom (DOF) than the task strictly requires, admit optimization-based IK formulations that generalize the null-space projection of Section~\ref{sec:sota-potential-fields} into a proper constrained program. Two architectural families matter for this thesis: a strictly prioritized cascade of QPs, in which avoidance is solved only within the null space left over by higher-priority tasks, and a single QP in which tracking, avoidance, and joint limits are all rows or cost terms of the same program --- the family this thesis's own controller belongs to.

Kanoun, Lamiraux, and Wieber represent the prioritized family, generalizing the classical task-priority framework to admit inequality tasks directly and solving a cascade of QPs that each respect the null space of every higher-priority task solved before it~\cite{kanoun2011kinematic}; Escande, Mansard, and Wieber's hierarchical quadratic programming formulates the same cascade as a single, efficiently factorized solve fast enough for online humanoid motion generation~\cite{escande2014hierarchical}. Within the single-QP family this thesis follows, two sub-lines differ in how avoidance enters the program. One treats avoidance as a cost: RelaxedIK poses every objective, including avoidance, as a weighted cost term solved at interactive rates~\cite{rakita2018relaxedik}, later extended with an explicit environment-distance penalty (CollisionIK)~\cite{rakita2021collisionik} and with ranged, rather than point, goals (RangedIK)~\cite{wang2023rangedik}; DawnIK applies the same cost-based principle to a decentralized, per-instant solve for heterogeneous multi-arm systems, at the cost of being myopic with no forward-invariance guarantee across successive solves~\cite{marangoz2023dawnik}. The other, architecturally closest to this thesis, treats avoidance as a hard constraint: Rauscher, Kimmel, and Hirche instantiate a CBF-QP directly at the joint-velocity level, the same level this thesis operates at, treating each link--obstacle pair as a barrier row~\cite{rauscher2016constrained}, and iKinQP, still a preprint, enforces minimum-distance inequalities directly as hard QP constraints in the same spirit~\cite{ashkanazy2023ikinqp}.

Three further pieces of infrastructure recur across this family and are reused directly in this thesis. Del Prete's analysis of joint bounds under discrete-time acceleration or torque control is the formal antecedent for a velocity-aware braking buffer rather than a fixed box~\cite{delprete2018joint}. At the trajectory level, Schulman et al.'s sequential convex optimization solves the same tracking-versus-avoidance trade-off over a whole trajectory by repeatedly convexifying a signed-distance cost~\cite{schulman2014trajopt}, and efficient, differentiable collision primitives~\cite{tracy2022diffpills,montaut2023diffcol} make gradient-based solvers of this kind tractable at interactive rates. All of the formulations above depend on a fast rigid-body kinematics backbone; the Pinocchio library provides exactly that layer, including the analytical derivatives needed for every Jacobian referenced here, and is the backbone this thesis's own controller is built on~\cite{carpentier2019}.

\subsection{Reference and command governors}
\label{sec:sota-governors}

A complementary strategy to constraining the control law itself is to constrain the reference it is asked to track before it reaches the controller. Garone, Di Cairano, and Kolmanovsky's survey of reference and command governors formalizes this as a minimally invasive outer layer: given a reference and a set of constraints the closed loop must respect, the governor computes the reference closest to the requested one that keeps the predicted trajectory admissible, modifying it only when the unmodified reference would cause a violation~\cite{garone2017}. Merckaert, Vanderborght, Nicotra, and Garone apply an explicit reference governor to a manipulator, filtering the Cartesian reference so the downstream tracking controller never receives a target that would force it outside joint or workspace limits, without adding any inequality row to the tracking controller's own optimization~\cite{merckaert2018constrained}. The governor and the CLF--CBF QP of Section~\ref{sec:sota-clf-cbf-qp} are complementary: the governor bounds what the controller is asked to do, the QP bounds what it is allowed to do in response, and neither substitutes for the other when the operator's raw command is both aggressive and close to an obstacle.

\subsection{Sampled-data, robust, and time-varying barrier functions}
\label{sec:sota-sampled-cbf}

The forward-invariance theorem of Section~\ref{sec:sota-cbf} assumes a continuously applied controller acting on an exactly known, continuous-time state; a QP solved once per tick on a sampled, possibly delayed measurement satisfies neither assumption exactly. A dedicated literature tightens the barrier condition to cover exactly this gap. Sampled-data CBF results shrink the enforced margin by a term proportional to the sampling period and to the barrier's own Lipschitz constant, so the discrete-time controller still certifies invariance despite acting only at isolated instants~\cite{breeden2021sampled}. A parallel line addresses disturbance rather than sampling: robust and input-to-state-safe CBFs bound how far the state may stray outside the nominal safe set under bounded model or actuation error~\cite{kolathaya2019issf}. Srinivasan, Santoyo, and Coogan's continuous transition between successive reachability tasks, already discussed in Section~\ref{sec:sota-positioning}, is the nearest instance of a changing constraint set actually cited in this thesis~\cite{srinivasan2020continuous}. \label{sec:sota-world-representation}
Measurement-robust CBFs~\cite{cosner2021mrcbf} address a related but distinct gap, widening the margin in proportion to perception uncertainty rather than to sampling or disturbance. None of these results is applied in this thesis, whose own barrier is stated and evaluated in nominal, continuous-time form over a sampled implementation.

\section{Perception for safety-critical control}
\label{sec:sota-perception}

This system does not re-estimate the world continuously: the head inspects the scene once, fits a collision model, and hands a frozen snapshot of it to the safety controller. Only two perception capabilities matter for that pipeline, and this section keeps to them: visually tracking the active arm, and fitting a collision model to an unknown tabletop from depth data.

\subsection{Visual tracking of the active arm}
\label{sec:sota-visual-servoing}

Keeping the active arm inside a head-mounted camera's field of view is itself a control problem. Roncone, Pattacini, Metta, and Natale's Cartesian gaze controller for a humanoid head resolves the competing demands of tracking a moving target while respecting the head's own joint limits and vergence constraint~\cite{roncone2016gaze}. Rakita, Mutlu, and Gleicher address an autonomous camera supporting a human teleoperator, showing that a camera policy anticipating where the operator's attention will be needed next, rather than only reacting to the manipulated object's current position, measurably improves remote task performance~\cite{rakita2018autonomouscamera}, a framing directly analogous to a head-mounted camera keeping a teleoperator's active hand framed while a shared-autonomy policy executes nearby.

\subsection{Building a collision model from tabletop perception}
\label{sec:sota-tabletop-world}

Converting a raw depth stream into obstacle geometry a safety-critical controller can consume is itself a substantial estimation pipeline. Fischler and Bolles's random sample consensus (RANSAC) algorithm, decades older than any application to robotics, remains the standard robust-fitting tool for extracting a dominant plane, a table surface, from a depth point cloud contaminated by outliers and sensor noise~\cite{fischler1981ransac}. Rusu, Marton, Blodow, Dolha, and Beetz's tabletop object-mapping pipeline is the closest prior architecture to a plane-then-cluster approach: it segments a table plane, clusters the points above it into candidate objects, and maintains the resulting object map incrementally, targeting exactly the household-manipulation tabletop scenario this thesis's own perception pipeline addresses~\cite{rusu2008tabletop}.

\section{Positioning of this thesis}
\label{sec:sota-positioning}

The preceding sections establish five propositions this thesis builds on: that a hard safety constraint dominates a soft penalty, both generally (Section~\ref{sec:sota-clf-cbf-qp}) and specifically after arbitration (Section~\ref{sec:sota-linear-blending-limits}); that arbitration keyed to directional agreement, rather than confidence or distance, better captures conflict between a human and an autonomous policy (Section~\ref{sec:sota-arbitration}); that a safety filter's own dual variables carry information worth surfacing beyond the controller computing them (Section~\ref{sec:sota-teleop-barriers}); that a barrier constraint is only as trustworthy as the perception pipeline supplying its geometry (Section~\ref{sec:sota-world-representation}); and that a forward-invariance guarantee is only as complete as the obstacle set it is stated over, which a manipulation task itself changes at the moments of grasp and release (Section~\ref{sec:sota-bimanual}). Five gaps emerge against the closest prior systems, from pure policy blending to the closest published post-blend safety filter, stated as precisely and conservatively as the literature search supports.

First, enforcing safety strictly after arbitration is established as necessary by Chaubey et al.'s formal argument and Guler et al.'s empirical one, but published realizations are approximate, position-space filters, at least among the systems identified in the literature reviewed here. BarrierIK solves a nonlinear program with a position-based update rule at each step, so its guarantee holds only to the tolerance of a local sequential least-squares quadratic programming (SLSQP) solve, validated in simulation and a virtual reality (VR) study rather than on hardware. This thesis instead applies the same post-blend principle inside a strictly convex QP formulated natively in joint velocity, with a first-order barrier condition, validated on a physical bimanual manipulator and evaluated systematically, at human-subject-study scale, in a matched simulation environment.

Second, arbitrating authority on the alignment between the operator's instantaneous twist and the policy's twist, rather than a scalar confidence, distance, or entropy summary, is barely explored, or at least not identified in the literature reviewed here. Oh, Toussaint, and Mainprice's disagreement-based arbitration is the nearest neighbour, trusting the operator more when independently reasonable strategies agree and less when they diverge, but it arbitrates on disagreement between sub-policies rather than alignment between the operator's own command and a single belief-weighted policy over a dynamically resolved goal manifold, and is not evaluated inside a safety-critical loop.

Third, using a safety-critical QP's own KKT dual variables, the shadow prices already computed at every solve, to adaptively schedule the QP's own constraint weights, rather than only as an external feedback channel, is, on the evidence surveyed here, unaddressed. Raphalen et al. is closest: it renders the multipliers as operator feedback but does not feed them back into the QP's own cost structure. The mechanism itself is implemented and exercised on hardware, not validated by a closed-loop stability analysis or an ablation against a fixed slack weight, a limitation this thesis states directly.

Fourth, no work identified in the literature reviewed here factorially compares both classical motion mappings of Section~\ref{sec:sota-mappings} with both assistance channels of Sections~\ref{sec:sota-shared-control} and~\ref{sec:sota-haptic-fixtures} in one study. Guler et al.'s study is closest by scale, six conditions and ten participants, but its conditions vary the safety mechanism on a single motion interface, VR controllers without force feedback, not the mapping itself.

Fifth, none of the systems identified in the literature reviewed here says what becomes of the barrier once an assisted grasp succeeds (Section~\ref{sec:sota-bimanual}). Manschitz et al. withhold their own collision-aware assistance for the entire held-object phase, calling it future work; Guler et al.'s safety formulation has no notion of a graspable object, so BarrierIK's pick-and-place evaluation never distinguishes an obstacle from something the operator has just picked up~\cite{manschitz2025samplinggrasp,guler2026barrierik}. The closest conceptual precedent lies outside manipulation entirely: Srinivasan, Santoyo, and Coogan introduce a time-varying term into a control-barrier-function constraint so a mobile robot can transition continuously between sequential reachability tasks rather than switching between quadratic programs discretely~\cite{srinivasan2020continuous}. Their transition is between successive navigation goals, never between an obstacle and a graspable object, and nothing in it addresses the continuity of the constraint row's own Lagrange multiplier through the change. This thesis instead treats a grasped object as a hybrid barrier: rather than a single barrier switched fully on or fully off, it passes through several distinct regimes as the grasp proceeds, each with its own rule for how the constraint is set. Only the safety margin is guaranteed to change smoothly within one regime; nothing is claimed about continuity across a change of regime.



The shared conclusion this thesis builds on is that enforcing safety after arbitration is what makes shared autonomy acceptable to an operator, and that soft penalty terms are no substitute for a hard constraint. Among the systems surveyed above, that conclusion has, until now, been demonstrated either without hardware or without haptic disclosure of the constraint, never alongside a constraint that schedules its own weight from its own shadow price. This thesis validates the controller and its perception pipeline on physical hardware, and separately evaluates the full factorial comparison at a simulated human-subject-study scale, $n=24$, no reviewed hardware study among those surveyed here reaching a comparable participant count, though the comparison itself is run in simulation rather than on hardware. The remaining chapters take up exactly that combination.


